\documentclass[a4paper,11pt]{article}

\usepackage[utf8]{inputenc}
\usepackage[T1]{fontenc}
\usepackage{geometry}
\geometry{margin=2.2cm}
\usepackage{xcolor}
\usepackage{graphicx}
\usepackage{tikz}
\usepackage{everypage}
\usepackage{background}
\usepackage{amsmath}
\usepackage{amssymb}
\usepackage{booktabs}
\usepackage{caption}
\usepackage{float}
\usepackage[colorlinks=true,linkcolor=blue,urlcolor=blue]{hyperref}

\setlength{\parskip}{0.6em}
\setlength{\parindent}{0pt}

\definecolor{bordercolor}{RGB}{68,94,114}
\definecolor{goodgreen}{RGB}{30,120,60}
\definecolor{missingred}{RGB}{180,40,40}

\backgroundsetup{
  scale=0.96, color=bordercolor, opacity=0.5, angle=0,
  position=current page.center,
  contents={
\begin{tikzpicture}[remember picture, overlay]
    \draw[line width=4pt, color=bordercolor]
      (current page.south west) rectangle (current page.north east);
  \end{tikzpicture}}
}

\begin{document}

\begin{titlepage}
\thispagestyle{empty}
\centering
\vspace*{3cm}
{\Huge\bfseries Probability of Default Scorecard\\\bigskip
\LARGE WoE-Based Retail PD Prototype with Out-of-Time Backtesting
\par}
\vspace{1.5cm}
{\Large Magnus Foldeström\\ June 21, 2026\par}
\vfill
\end{titlepage}

\section{What this project does}

This is a retail Probability of Default (PD) scorecard prototype built
end-to-end on a synthetic loan portfolio. The model follows the
methodology used by European banks for IRB scorecards: Weight-of-Evidence
(WoE) binning, Information-Value-based feature selection, regularized
logistic regression on WoE-encoded features, intercept calibration to
the long-run-average default rate, mapping to a rating master scale,
and out-of-time backtesting.

It is a portfolio demonstration, not a regulator-approved model. The
goal is to show that I can implement the structural components of an
IRB scorecard end-to-end and reason about how they connect to Basel
and IFRS 9 requirements. A production-grade version would need real
bureau data, a defined default trigger, independent validation, and
the full model governance infrastructure described in the limitations
section.

\section{Data}

A synthetic portfolio of 250,000 unsecured retail loans, 25,000 per
origination vintage across ten half-year cohorts (2019H1 through
2023H2). The dataset spans five years, which matches the Basel minimum
observation period for retail IRB PD estimation. Each loan has an
18-month performance window and a realized default outcome.

The data-generating process embeds a logit model where bureau
attributes (score, prior bankruptcies, worst delinquency status,
inquiries, utilization) are the strong predictors, application
attributes (DTI, income, employment) are moderate predictors, and
demographics (age, loan purpose) are weak predictors. Vintage-level
fixed effects produce cohort-to-cohort variation in the realized
default rate, mimicking the macroeconomic path from the pre-COVID
benign cohorts, through the 2020 stress (masked by forbearance) and
the 2021 recovery, and into the 2022--2023 rate-hike pressure.

\begin{table}[H]
\centering
\begin{tabular}{l r r r}
\toprule
\textbf{Vintage} & \textbf{Loans} & \textbf{Defaults} & \textbf{Default rate} \\
\midrule
2019H1 & 25\,000 & 2\,513 & 10.05\% \\
2019H2 & 25\,000 & 2\,556 & 10.22\% \\
2020H1 & 25\,000 & 2\,750 & 11.00\% \\
2020H2 & 25\,000 & 2\,939 & 11.76\% \\
2021H1 & 25\,000 & 2\,866 & 11.46\% \\
2021H2 & 25\,000 & 2\,764 & 11.06\% \\
2022H1 & 25\,000 & 3\,041 & 12.16\% \\
2022H2 & 25\,000 & 3\,141 & 12.56\% \\
2023H1 & 25\,000 & 3\,399 & 13.60\% \\
2023H2 & 25\,000 & 3\,605 & 14.42\% \\
\midrule
Total  & 250\,000 & 29\,574 & 11.83\% \\
\bottomrule
\end{tabular}
\caption{Portfolio composition by origination vintage. Realized default rates trend from 10.05\% in the pre-COVID cohort to 14.42\% in the 2023H2 stress cohort.}
\end{table}

\subsection{Cohort split}

The split is chronological by vintage, as is standard in IRB modelling,
because covariate drift across cohorts is the dominant failure mode in
production:

\begin{itemize}
\item \textbf{Training:} 2019H1 through 2021H2 (six cohorts, 150\,000 loans).
\item \textbf{Validation:} 2022H1 + 2022H2 (two cohorts, 50\,000 loans), in-time.
\item \textbf{Out-of-time:} 2023H1 + 2023H2 (two cohorts, 50\,000 loans), the
worst-performing cohorts in the panel, used as a stress test.
\end{itemize}

\section{Methodology}

\subsection{Weight of Evidence binning}

Each feature is transformed from raw value to a WoE score that
captures its monotonic relationship with default. For a bin $b$ of
feature $X$, WoE is defined as
$$\mathrm{WoE}_b = \log\!\frac{\Pr(X \in b \mid y = 0)}{\Pr(X \in b \mid y = 1)}$$
positive when non-defaulters are over-represented in the bin, negative
otherwise.

Numeric features are binned into deciles. Low-cardinality numeric
features (counts of bankruptcies, public records, missingness flags)
use their distinct values directly as bins. Categorical features
(home ownership, loan purpose) are encoded with WoE per category.
All encoders are fit on training data only and applied unchanged to
validation and OOT cohorts.

\subsection{Information Value and feature selection}

For each feature, the Information Value summarizes how much the
binned distribution of $X$ differs between defaulters and
non-defaulters:
$$\mathrm{IV} = \sum_b \big( \Pr(X \in b \mid y = 0) - \Pr(X \in b \mid y = 1) \big) \cdot \mathrm{WoE}_b$$

The industry convention is to drop features with IV below 0.02, on
the basis that anything weaker is essentially unpredictive and only
adds noise to the regression.

\begin{table}[H]
\centering
\begin{tabular}{l r}
\toprule
\textbf{Feature} & \textbf{IV} \\
\midrule
bureau\_score                          & 0.601 \\
interest\_rate\_pct                    & 0.512 \\
utilization\_pct                       & 0.476 \\
worst\_status\_24m                     & 0.331 \\
delinquencies\_2y                      & 0.283 \\
months\_since\_last\_delinquency\_missing & 0.263 \\
months\_since\_last\_delinquency       & 0.140 \\
prior\_bankruptcies                    & 0.105 \\
inquiries\_6m                          & 0.097 \\
public\_records                        & 0.026 \\
\midrule
\textit{Below 0.02 threshold (dropped)} & \textit{16 features} \\
\bottomrule
\end{tabular}
\caption{Information Value per feature, kept set above the 0.02 threshold. Bureau-derived attributes dominate as expected.}
\end{table}

\begin{figure}[H]
\centering
\includegraphics[width=0.85\textwidth]{02_information_values.png}
\caption{Information Value ranking. Bars above the red threshold are kept; below are dropped.}
\end{figure}

The final model uses ten features. All are bureau-derived or strongly
bureau-influenced (interest rate, set by the lender based on the
applicant's risk profile). Application-only attributes (DTI,
employment length, income, age) all fall below the IV threshold once
the bureau-based features are in.

\subsection{Logistic regression on WoE features}

The model is L2-regularized logistic regression with $C = 1.0$:
$$\log\!\frac{p_i}{1 - p_i} = \beta_0 + \sum_j \beta_j \cdot \mathrm{WoE}_j(x_{ij})$$
fitted by maximum likelihood on the training cohort.

\subsection{Calibration to long-run-average PD}

The raw logistic output is rank-discriminative but the level can drift
when training and prediction populations differ. The standard IRB fix
is an additive intercept correction:
$$\widehat{\mathrm{PD}}_i = \frac{1}{1 + \exp\!\big( -(\mathrm{logit}(\widehat p_i^{\text{raw}}) + c) \big)}$$
where $c$ is chosen so that the mean of $\widehat{\mathrm{PD}}$ on the
training set equals the realized default rate across the six training
cohorts. Because the training panel spans pre-COVID, COVID, and
recovery, the resulting calibration target approximates a
through-the-cycle average rather than any single vintage.

\subsection{Master rating scale}

PDs are mapped to a nine-grade scale (A1 to C3), constructed as
piecewise cutpoints in PD space. The grades are the unit of currency
for the rest of the bank: loan pricing, capital calculation, and IFRS
9 staging all run on grades, not on raw PD scores. The grade
boundaries are 0.5\%, 1\%, 2.5\%, 5\%, 8\%, 12\%, 18\%, 28\%, and
100\%.

\section{Results}

\subsection{Discrimination}

\begin{table}[H]
\centering
\begin{tabular}{l r r r r r}
\toprule
\textbf{Cohort} & \textbf{AUC} & \textbf{Gini} & \textbf{KS} & \textbf{Log-loss} & \textbf{Brier} \\
\midrule
TRAIN (150k, 2019--2021) & 0.7778 & 0.5557 & 0.4126 & 0.2915 & 0.0846 \\
VAL 2022 (50k)            & 0.7699 & 0.5397 & 0.4038 & 0.3210 & 0.0951 \\
OOT 2023 (50k)            & 0.7693 & 0.5386 & 0.3991 & 0.3492 & 0.1048 \\
\bottomrule
\end{tabular}
\caption{Headline metrics across the chronological split. Gini sits in the upper half of the typical retail PD range (0.40--0.60 acceptable, 0.55+ strong) and drops only 1.7 Gini-points from train to out-of-time despite spanning a full macroeconomic cycle.}
\end{table}

\begin{figure}[H]
\centering
\includegraphics[width=0.78\textwidth]{01_roc_three_splits.png}
\caption{ROC curves across train, validation, and out-of-time. The three curves are nearly superimposable.}
\end{figure}

\subsection{Calibration per rating grade}

\begin{table}[H]
\centering
\footnotesize
\begin{tabular}{l r r r r}
\toprule
\textbf{Grade} & \textbf{N (OOT)} & \textbf{Pred PD} & \textbf{Observed DR} & \textbf{Share} \\
\midrule
A2  & 1\,314 & 0.86\%  & 0.91\%  & 2.6\% \\
A3  & 7\,136 & 1.77\%  & 2.59\%  & 14.3\% \\
B1  & 10\,675 & 3.68\%  & 5.56\%  & 21.4\% \\
B2  & 8\,413 & 6.37\%  & 9.14\%  & 16.8\% \\
B3  & 7\,202 & 9.84\%  & 12.77\% & 14.4\% \\
C1  & 6\,160 & 14.75\% & 18.86\% & 12.3\% \\
C2  & 5\,041 & 22.28\% & 28.47\% & 10.1\% \\
C3  & 4\,059 & 39.95\% & 47.50\% & 8.1\% \\
\bottomrule
\end{tabular}
\caption{Predicted vs observed default rate per grade on the 2023 out-of-time cohort. The TTC-calibrated model under-predicts realized rates in the B and C grades by 2--7 percentage points, reflecting the macroeconomic stress that a TTC scorecard is not designed to track point-in-time.}
\end{table}

\begin{figure}[H]
\centering
\includegraphics[width=0.92\textwidth]{04_grade_calibration_oot.png}
\caption{Predicted PD vs observed default rate per grade on the 2023 out-of-time cohort.}
\end{figure}

\begin{figure}[H]
\centering
\includegraphics[width=0.78\textwidth]{03_calibration_oot.png}
\caption{Calibration curve on OOT, 10 quantile bins. The systematic upward shift across the range is the expected behaviour of a TTC scorecard during a downturn.}
\end{figure}

\subsection{Population stability}

The Population Stability Index measures whether the grade distribution
on a new cohort matches the training distribution:
$$\mathrm{PSI} = \sum_g (\pi_g^{\text{new}} - \pi_g^{\text{train}}) \cdot \log\!\frac{\pi_g^{\text{new}}}{\pi_g^{\text{train}}}$$

\begin{table}[H]
\centering
\begin{tabular}{l r}
\toprule
\textbf{Comparison} & \textbf{PSI} \\
\midrule
TRAIN vs VAL (2022)  & 0.0003 \\
TRAIN vs OOT (2023)  & 0.0001 \\
\bottomrule
\end{tabular}
\caption{PSI well below the conventional thresholds (0.10 stable, 0.25 material shift). The grade distribution barely moves across the five-year panel.}
\end{table}

\begin{figure}[H]
\centering
\includegraphics[width=0.92\textwidth]{05_population_stability.png}
\caption{Share of population in each grade across train, validation, and OOT cohorts.}
\end{figure}

\section{Regulatory framework}

\subsection{Basel II and III}

Basel II (2004) introduced the Internal Ratings-Based (IRB) approach,
which lets banks use their own PD, LGD, and EAD estimates to compute
regulatory capital. The IRB capital formula for a single retail
exposure depends most heavily on PD: a higher PD drives a higher risk
weight, which drives a higher required capital reserve. Basel III
(2010 onwards) tightened the capital requirements, added the leverage
ratio and liquidity coverage ratio, and constrained how much banks
could rely on internal models versus standardised risk weights.

\subsection{IFRS 9}

IFRS 9, effective from January 2018, replaced the incurred-loss model
of IAS 39 with the Expected Credit Loss (ECL) framework for loan
provisioning. Every loan is assigned to one of three stages:

\begin{itemize}
\item \textbf{Stage 1}: performing. ECL provision equals the expected
loss over the next 12 months.
\item \textbf{Stage 2}: significant increase in credit risk since
origination. ECL provision equals the expected loss over the loan's
remaining life.
\item \textbf{Stage 3}: credit-impaired (default). Lifetime ECL on a
defaulted basis.
\end{itemize}

The PD that drives staging and ECL under IFRS 9 differs from Basel in
two ways. IFRS 9 uses a point-in-time (PIT) PD conditional on current
macroeconomic conditions, and requires lifetime PD for Stage 2 and 3.

\section{Compliance checklist}

The tables below list the specific Basel and IFRS 9 requirements
relevant to a retail PD scorecard, marked against what is present in
this prototype and what is not.

\subsection{Basel II/III requirements}

\begin{table}[H]
\centering
\footnotesize
\begin{tabular}{l c l}
\toprule
\textbf{Requirement} & \textbf{Status} & \textbf{In this project} \\
\midrule
Minimum 5 years of observation data & \textcolor{goodgreen}{\textbf{Met}} & 2019--2023, ten half-year vintages \\
Master scale, $\geq$ 7 non-default grades & \textcolor{goodgreen}{\textbf{Met}} & 9-grade A1--C3 scale \\
Through-the-cycle PD calibration & \textcolor{goodgreen}{\textbf{Met}} & Intercept-shifted to long-run-average \\
Chronological out-of-time validation & \textcolor{goodgreen}{\textbf{Met}} & 2023 cohorts held out for backtest \\
Discrimination metrics reported & \textcolor{goodgreen}{\textbf{Met}} & Gini, KS, AUC per cohort \\
Stability monitoring on rating output & \textcolor{goodgreen}{\textbf{Met}} & PSI on grade distribution \\
Reasonable sample per grade & \textcolor{goodgreen}{\textbf{Met}} & Thousands of defaults per grade \\
\midrule
12-month PD horizon as regulatory anchor & \textcolor{missingred}{\textbf{Partial}} & 18-month performance window used \\
\midrule
Defined default trigger (90+ DPD or UTP) & \textcolor{missingred}{\textbf{Not in scope}} & Binary flag, no DPD/UTP construct \\
Override capability & \textcolor{missingred}{\textbf{Not in scope}} & No manual override layer built \\
Independent validation & \textcolor{missingred}{\textbf{Not in scope}} & Prototype, no second-line review \\
Use test (real lending decisions) & \textcolor{missingred}{\textbf{Not in scope}} & Demonstration only \\
Model risk management framework & \textcolor{missingred}{\textbf{Not in scope}} & No MRM infrastructure \\
Supervisory approval (Finansinspektionen) & \textcolor{missingred}{\textbf{Not in scope}} & Portfolio project, not submitted \\
\bottomrule
\end{tabular}
\caption{Basel II/III compliance checklist. The structural and quantitative requirements relevant to model development are met. The governance and process requirements that a production deployment requires are out of scope for a portfolio prototype.}
\end{table}

\subsection{IFRS 9 requirements}

\begin{table}[H]
\centering
\footnotesize
\begin{tabular}{l c l}
\toprule
\textbf{Requirement} & \textbf{Status} & \textbf{In this project} \\
\midrule
Calibrated PD foundation suitable for ECL & \textcolor{goodgreen}{\textbf{Met}} & TTC-calibrated PD at grade level \\
Origination-grade infrastructure for SICR & \textcolor{goodgreen}{\textbf{Met}} & Master scale enables grade comparison \\
\midrule
Stage 1/2/3 classification implementation & \textcolor{missingred}{\textbf{Not built}} & Infrastructure present, logic not coded \\
SICR triggers (PD ratio, days past due) & \textcolor{missingred}{\textbf{Not built}} & No SICR rule implemented \\
Lifetime PD projection & \textcolor{missingred}{\textbf{Not built}} & Single-horizon PD only \\
PIT PD via macro overlay & \textcolor{missingred}{\textbf{Not built}} & Only TTC version computed \\
Multi-scenario weighting (base/up/down) & \textcolor{missingred}{\textbf{Not built}} & Single deterministic scenario \\
Discounted lifetime ECL calculation & \textcolor{missingred}{\textbf{Not built}} & ECL formula not applied \\
LGD model & \textcolor{missingred}{\textbf{Not built}} & Out of scope for PD project \\
EAD model & \textcolor{missingred}{\textbf{Not built}} & Out of scope for PD project \\
\bottomrule
\end{tabular}
\caption{IFRS 9 compliance checklist. The PD model provides the foundation but the ECL engine, staging logic, macro overlay, and LGD/EAD models that complete IFRS 9 compliance are scheduled as separate projects in the portfolio.}
\end{table}

\subsection{Summary}

This prototype covers the methodological core of a Basel IRB PD
scorecard. The gaps are the governance and infrastructure layers that
sit above the model (independent validation, MRM framework, use test,
supervisory approval) and the IFRS 9 ECL engine that sits below it
(staging, lifetime PD, macro overlay, LGD, EAD). Both gaps are
expected for a portfolio project rather than a production deployment.

\section{What the model shows}

The scorecard discriminates reasonably (Gini $\approx 0.54$ on OOT),
is tightly calibrated per grade on the training data, and shows the
expected mild under-prediction during the 2023 stress vintage. PSI
confirms that the grade distribution is stable across the full
five-year panel.

The headline pattern is the OOT under-prediction. Observed default
rates exceed predicted by 2--7 percentage points in the higher-risk
grades. This is the textbook signature of a TTC-calibrated scorecard
during a downturn: the model is anchored to a long-run-average that
includes the benign pre-COVID years, and the 2023 cohort runs above
that average. In production this would be absorbed by a PIT macro
overlay for IFRS 9 ECL while keeping the TTC version for Basel
regulatory capital.

\section{Limitations}

\begin{itemize}
\item Synthetic data. The bureau attributes are simulated, not from a
real credit bureau. The IV pattern is plausible but not validated
against external benchmarks.
\item No defined default trigger. The target is a binary outcome flag
without 90+ days past due or unlikely-to-pay construction.
\item Single performance window. The 18-month outcome horizon is
fixed across vintages.
\item No monotone constraints on WoE. Some features end up with
non-monotone WoE curves, which is statistically fine but a model
governance committee would typically require monotonicity.
\item No fairness audit. A production model would test for disparate
impact across protected attributes.
\item No champion-challenger framework. A production deployment would
shadow the new model against the existing scorecard for 3--6 months
before any traffic switch.
\item Governance and infrastructure layers (independent validation,
MRM framework, use test, supervisory approval) are out of scope.
\end{itemize}

\section{References}

\begin{itemize}
\item Basel Committee on Banking Supervision, \emph{International
Convergence of Capital Measurement and Capital Standards} (Basel II/III).
\item European Banking Authority, \emph{Guidelines on PD Estimation,
LGD Estimation and the Treatment of Defaulted Exposures} (EBA/GL/2017/16).
\item Siddiqi, \emph{Credit Risk Scorecards: Developing and Implementing
Intelligent Credit Scoring} (Wiley, 2012). Standard industry reference
for the WoE/IV scorecard workflow.
\item Anderson, \emph{The Credit Scoring Toolkit} (Oxford, 2007).
\item IFRS 9, \emph{Financial Instruments}. The accounting standard
that drives the staging engine and lifetime ECL.
\end{itemize}

\end{document}
